%% Font size %%
\documentclass[11pt]{article}

%% Load the custom package
\usepackage{Mathdoc}

%% Numéro de séquence %% Titre de la séquence %%
\renewcommand{\centerhead}{Décimaux - Éval 1 : Décomposer}

%% Spacing commands %%
\renewcommand{\baselinestretch}{1} \setlength{\parindent}{0pt}

\begin{document}

\phantom{0}
%\vspace{-1.5cm}

\entetedevoirs{25}

\begin{exercicedevoir}[5][Compléter les phrases suivantes]
\textit{Exemple :} Le chiffre des dizaines du nombre $53{,}6$ est : 5
\begin{enumerate}
\item Le chiffre des unités du nombre $3\,927{,}801$ est : $\ldots\ldots\ldots\ldots$
\item Le chiffre des millièmes du nombre $3\,927{,}801$ est : $\ldots\ldots\ldots\ldots$
\item Le chiffre des centièmes du nombre $3\,927{,}801$ est : $\ldots\ldots\ldots\ldots$
\item Le chiffre des dixièmes du nombre $3\,927{,}801$ est : $\ldots\ldots\ldots\ldots$
\item Le chiffre des centaines du nombre $3\,927{,}801$ est : $\ldots\ldots\ldots\ldots$
\end{enumerate}
\end{exercicedevoir}

\begin{exercicedevoir}[5][Donner les formes décimales I]
\textit{Exemple :} Forme décimale de $7 + \dfrac{56}{100} =7,56$
\begin{enumerate}
\item Forme décimale de $5050 + \dfrac{6}{1000} + \dfrac{7}{10} + \dfrac{7}{100}=\ldots\ldots\ldots\ldots$
\item Forme décimale de $386 + \dfrac{7}{100}=\ldots\ldots\ldots\ldots$
\item Forme décimale de $71 + \dfrac{4}{100} + \dfrac{2}{10} + \dfrac{6}{1000}=\ldots\ldots\ldots\ldots$
\item Forme décimale de $5 + \dfrac{7}{1000}=\ldots\ldots\ldots\ldots$
\item Forme décimale de $17 + \dfrac{7}{100} + \dfrac{8}{10} + \dfrac{4}{1000}=\ldots\ldots\ldots\ldots$
\end{enumerate}
\end{exercicedevoir}

\begin{exercicedevoir}[5][Donner les formes décimales II]
\textit{Exemple :} Forme décimale de $7 + \dfrac{56}{100} =7,56$
\begin{enumerate}
\item Forme décimale de $7160 + \dfrac{5}{100} =\ldots\ldots\ldots\ldots$
\item Forme décimale de $653 + \dfrac{3}{10}=\ldots\ldots\ldots\ldots$
\item Forme décimale de $12 + \dfrac{71}{1000}=\ldots\ldots\ldots\ldots$
\item Forme décimale de $929 + \dfrac{2}{100}=\ldots\ldots\ldots\ldots$
\item Forme décimale de $118 + \dfrac{381}{1000}=\ldots\ldots\ldots\ldots$
\end{enumerate}
\end{exercicedevoir}

\begin{exercicedevoir}[5][Écrire sous la forme décomposée I]
\textit{Exemple :} $7,56 = 7 + \dfrac{5}{10} +\dfrac{6}{100}$
\begin{enumerate}
\item $34,572=\ldots\ldots\ldots\ldots+\dfrac{\ldots\ldots\ldots\ldots}{\ldots\ldots\ldots\ldots}+\dfrac{\ldots\ldots\ldots\ldots}{\ldots\ldots\ldots\ldots}+\dfrac{\ldots\ldots\ldots\ldots}{\ldots\ldots\ldots\ldots}$
\item $8,407=\ldots\ldots\ldots\ldots+\dfrac{\ldots\ldots\ldots\ldots}{\ldots\ldots\ldots\ldots}+\dfrac{\ldots\ldots\ldots\ldots}{\ldots\ldots\ldots\ldots}+\dfrac{\ldots\ldots\ldots\ldots}{\ldots\ldots\ldots\ldots}$
\item $62,310=\ldots\ldots\ldots\ldots+\dfrac{\ldots\ldots\ldots\ldots}{\ldots\ldots\ldots\ldots}+\dfrac{\ldots\ldots\ldots\ldots}{\ldots\ldots\ldots\ldots}+\dfrac{\ldots\ldots\ldots\ldots}{\ldots\ldots\ldots\ldots}$
\item $15,006=\ldots\ldots\ldots\ldots+\dfrac{\ldots\ldots\ldots\ldots}{\ldots\ldots\ldots\ldots}+\dfrac{\ldots\ldots\ldots\ldots}{\ldots\ldots\ldots\ldots}+\dfrac{\ldots\ldots\ldots\ldots}{\ldots\ldots\ldots\ldots}$
\item $90,273=\ldots\ldots\ldots\ldots+\dfrac{\ldots\ldots\ldots\ldots}{\ldots\ldots\ldots\ldots}+\dfrac{\ldots\ldots\ldots\ldots}{\ldots\ldots\ldots\ldots}+\dfrac{\ldots\ldots\ldots\ldots}{\ldots\ldots\ldots\ldots}$
\end{enumerate}
\end{exercicedevoir}

\begin{exercicedevoir}[5][Écrire sous la forme décomposée II]
\textit{Exemple :} $7,56 = 7 + \dfrac{56}{100}$
\begin{enumerate}
\item $45,83=\ldots\ldots\ldots\ldots+\dfrac{\ldots\ldots\ldots\ldots}{\ldots\ldots\ldots\ldots}$
\item $12,407=\ldots\ldots\ldots\ldots+\dfrac{\ldots\ldots\ldots\ldots}{\ldots\ldots\ldots\ldots}$
\item $0,5=\ldots\ldots\ldots\ldots+\dfrac{\ldots\ldots\ldots\ldots}{\ldots\ldots\ldots\ldots}$
\item $6,284=\ldots\ldots\ldots\ldots+\dfrac{\ldots\ldots\ldots\ldots}{\ldots\ldots\ldots\ldots}$
\item $99,0682=\ldots\ldots\ldots\ldots+\dfrac{\ldots\ldots\ldots\ldots}{\ldots\ldots\ldots\ldots}$
\end{enumerate}
\end{exercicedevoir}


\end{document}

%%% Local Variables:
%%% mode: LaTeX
%%% TeX-master: t
%%% TeX-master: t
%%% End:

